\documentclass[11pt]{report}
% T0_preamble_standalone_A4_En.tex — LOKALER PLATZHALTER
% Im Repo durch die offizielle Korpus-Präambel ersetzen; die
% \input-Zeilen der Wrapper bleiben unverändert gültig.
\usepackage{fontspec}
\usepackage[english]{babel}
\usepackage{csquotes}
\setmainfont{Inter}[BoldFont=Inter-Bold,ItalicFont=Inter-Italic,BoldItalicFont=Inter-BoldItalic]
\setmonofont{JetBrains Mono}[BoldFont=JetBrainsMono-Bold,ItalicFont=JetBrainsMono-Italic]
\usepackage{unicode-math}
\setmathfont{Libertinus Math}
\usepackage{geometry}
\geometry{a4paper,top=2.4cm,bottom=2.4cm,left=2.8cm,right=2.8cm}
\usepackage{microtype,parskip,xcolor}
\usepackage{tcolorbox}
\tcbuselibrary{breakable,skins}
\usepackage{booktabs,longtable,array,amsmath,physics,siunitx}
\sisetup{locale=US}
\usepackage{hyperref}
\hypersetup{colorlinks=true,linkcolor=blue!60!black,urlcolor=green!40!black}

% T0_preamble_patches.tex — LOKALER PLATZHALTER
% Im Repo durch die offizielle Version ersetzen.


\title{Doc.~311 -- Four onto Three\\
\large How does a 4D lattice arrive at three spatial dimensions?}
\author{Johann Pascher}
\date{1 August 2026}

\begin{document}
\maketitle
\tableofcontents

\chapter*{Doc.~311 -- Four onto Three}
\addcontentsline{toc}{chapter}{Doc.~311 -- Four onto Three}

\section*{What this is about}

A lattice in $\mathbb{R}^4$ has four directions on an equal footing.
Measurement takes place in three. How does one get from the four to the
three?

There are exactly three routes, and their costs differ greatly. This
note walks through them, condenses them at the end into two conditions,
and applies those conditions to three cases: the author's own
construction, generic projection constructions, and string theory.

\medskip
\noindent\textbf{Status:} working note. Not part of the A-Series. The
chapters on the routes and on the decomposition finding are computed
(\texttt{ffgft\_311\_vier\_auf\_drei.py}); the chapter
\enquote{Application} poses a question and makes no assertion.

\chapter{The three routes}

\section{Route 1 --- Compactify}

The fourth direction becomes a circle of radius $R$. Kaluza 1921,
Klein 1926.

\textbf{What it buys.} On a circle only integer windings are admitted,
and from that a mode tower arises with $m_n \propto n/R$. The mode
index \emph{is} a scale. Nothing has to be attached.

\textbf{What it costs.} $R$ is a free modulus. Whoever does not say
what fixes $R$ has gained a scale and acquired a parameter --- that is,
nothing. This is the price string theory pays with its landscape.

The modulus disappears only when $R$ is bound to another quantity. In
FFGFT this happens through $\tilde T \cdot m = 1$: once the mass is
fixed, the cycle is fixed, and nothing is left to tune.

\section{Route 2 --- Project}

The 4D lattice is projected onto a 3D section. The standard route for
quasicrystals, cut-and-project.

\textbf{What it costs.} The projection has a window parameter, and the
lattice invariants do not survive it. For D4 onto the first three
coordinates:
\[
  24 \text{ shortest vectors} \;\longrightarrow\; 18 \text{ distinct images.}
\]
Six pairs coincide. What arrives after projection is no longer a
lattice with kissing number~24, but an aperiodic point pattern with
different invariants.

\section{Route 3 --- The fourth is not space}

The fourth coordinate carries a different quantity: time, mass, charge,
an internal index. Then only three directions are space, and there is
nothing to reduce.

\textbf{What it costs.} One must say \emph{what} it carries, and the
statement must bear weight. This is no formality: a fourth direction
called \enquote{mass} without anything binding it to a mass is merely a
label.

\chapter{Routes 1 and 3: each incomplete on its own}

The two look alike because both take the fourth direction out of space.
But they do different things, and the difference is precisely the point
at which a construction holds or does not.

\textbf{Route 1 makes the direction compact and leaves open what it
carries.} Pure Kaluza-Klein. A circle of radius $R$ enforces integer
windings, from that a mode tower $m_n \propto n/R$ arises, and the mode
index is a scale. But $R$ itself is free. One has gained a scale and
bought a parameter along with it.

\textbf{Route 3 says what the direction carries and leaves open whether
it is compact.} It is called \enquote{mass} or \enquote{charge} or
\enquote{internal index}. Without compactness this is a label: nothing
enforces integer values, without integrality there is no mode tower,
and without a mode tower no scale.

\textbf{Each is incomplete on its own, in opposite ways:}
\begin{center}
\begin{tabular}{ll}
Route 1: & mechanism without meaning \\
Route 3: & meaning without mechanism \\
\end{tabular}
\end{center}

\section{What closes the gap}

Both at once, and the binding between them.

\begin{center}
\begin{tabular}{lllp{4.6cm}}
\toprule
 & compact & bound & result \\
\midrule
Kaluza-Klein & yes & no & scale present, but free modulus \\
mere label & no & in name only & no mode tower, no scale \\
\textbf{FFGFT} & \textbf{yes} & \textbf{yes, via $\tilde T \cdot m = 1$}
  & \textbf{scale without modulus} \\
\bottomrule
\end{tabular}
\end{center}

The reciprocity closes both gaps in one step. The mode tower acquires a
meaning because the direction is the mass direction. And the free
radius disappears because $R$ is no longer choosable once $m$ is fixed
--- the cycle \emph{is} the mass, not a container for it.

\section{The sharper question}

The real distinction is therefore not \enquote{which of the three
routes}, but:

\begin{quote}
Is the fourth direction compact \textbf{and} bound to something --- or
only one of the two?
\end{quote}

Whoever only compactifies has a modulus to stabilise. Whoever only
names has a word. Both together and bound to each other is the only
case in which a scale arises without a parameter following behind.

This is also the touchstone for other people's constructions, and it
can be put in a single question: \emph{what fixes the radius?} If no
answer comes, the scale has merely been displaced.

\chapter{Route 2 in detail: what projecting presupposes and costs}

Route 2 differs from the other two not in degree but in kind.

\section{Projection needs a standpoint outside}

Compactifying is a property of the object: the circle is there or it is
not, regardless of who is looking. Interpreting the fourth direction as
the mass direction is a statement about the structure.

Projecting is a map \textbf{from outside}. It needs a cut direction and
a window, and both are chosen by the observer. They say nothing about
the lattice; they say something about the access to it.

This is a different type of free parameter from the radius in Route~1.
A radius is a feature of the object and can be bound to another
quantity of the object --- exactly what $\tilde T \cdot m = 1$
achieves. A cut direction cannot be bound. It can only be fixed, and
the fixing remains a stipulation.

\section{What happens to the 24 vectors}

Under projection onto the first three coordinates:
\begin{center}
\begin{tabular}{lcl}
24 vectors & $\longrightarrow$ & 18 images \\[4pt]
12 FCC vectors ($x_4 = 0$) & $\longrightarrow$
  & 12 images, length $\sqrt2$, untouched \\
12 winding vectors & $\longrightarrow$
  & \phantom{1}6 images, length 1, merged in pairs \\
\end{tabular}
\end{center}

For $(1,0,0,+1)$ and $(1,0,0,-1)$ have the same image $(1,0,0)$. The
difference between two opposite windings disappears.

And the image contains two different lengths, $1$ and $\sqrt2$. The
image is no longer a lattice with a uniform kissing number; packing
density in the usual sense loses its meaning there.

\textbf{This is the real loss.} Under compactification the twelve
winding neighbours remain distinguishable --- $+1$ and $-1$ are
different windings and carry different quantum numbers. Under
projection they are made indistinguishable. What leads out of space is
erased rather than reinterpreted.

\section{Where the irrational numbers come from}

For a rational cut direction, say a coordinate axis, finitely many
image lengths remain --- here two. The shortening of the winding
vectors by $1/\sqrt2$ is irrational, but it repeats.

In genuine cut-and-project the cut direction is chosen irrational, and
then the projection lengths are generically irrational. For the
direction $(1, \varphi, 0)$:
\begin{center}
\begin{tabular}{lcr}
$(1, 0, 0)$ & $\longrightarrow$ & $0.525731$ \\
$(0, 1, 0)$ & $\longrightarrow$ & $0.850651$ \\
$(1, 1, 0)$ & $\longrightarrow$ & $1.376382$ \\
$(1,-1, 0)$ & $\longrightarrow$ & $-0.324920$ \\
\end{tabular}
\end{center}

Every length a new number, and the image is \textbf{dense}. Only the
window makes it discrete again --- the second free parameter.

This is exactly how quasicrystals arise. And exactly how $\varphi$
enters there:

\begin{quote}
The stretching factors are irrational, yes. But they are properties of
the \textbf{chosen direction}, not of the lattice. Choose differently
and they are different numbers.
\end{quote}

\section{The consequence for other constructions}

Whoever brings $\varphi$ in through a cut-and-project construction has
put it in. It does not come from the lattice but from the choice of
projection direction. That is admissible, provided it is said --- and
it is something fundamentally different from a number that follows from
the structure.

The touchstone from the previous chapter thus acquires a second
question beside \emph{what fixes the radius?}:

\begin{quote}
\textbf{Who chose the cut direction, and from what?}
\end{quote}

If no answer comes, all numbers originating from the projection are
properties of the observer.

\chapter{The reversal: where \texorpdfstring{$\varphi$}{phi} is found,
the direction is no longer free}

If $\varphi$ appears in cut-and-project because the cut direction was
chosen accordingly --- what holds when $\varphi$ is not chosen but
\emph{found}?

\section{Where \texorpdfstring{$\varphi$}{phi} stands in the corpus}

A120 derives $\theta = 2/9$ from the redistribution of the lepton modes
under a five-fold rotation. The three weights:
\begin{center}
\begin{tabular}{llll}
$p_0 = 2/9$ & $= 0.2222222$ & \textbf{rational} & --- no $\varphi$ in it \\
$p_2 = (5 - 3\varphi)/9$ & $= 0.0162109$ & contains $\varphi$ & \\
$p_1 = (2 + 3\varphi)/9$ & $= 0.7615669$ & contains $\varphi$ & \\
\cmidrule{2-2}
sum & $= 1.0000000$ & & \\
\end{tabular}
\end{center}

Remarkably, the share passing into the trivial mode is
\textbf{rational}. $\varphi$ sits solely in the split between the two
non-trivial modes. All three carry the denominator $9 = 3^2$, the same
Z$_3$ structure in which $2/9$ sits as an orbit address.

\section{\texorpdfstring{$\varphi$}{phi} is not a choice here but an
eigenvalue}

In cut-and-project $\varphi$ appears because the observer laid the
direction that way. Here $\varphi$ is the eigenvalue of the structure
itself:
\[
  \varphi^2 = \varphi + 1 .
\]
A$_5$ is the icosahedral group. Its five-fold axes are the twelve
vertices with coordinates $(0, \pm1, \pm\varphi)$ and cyclic
permutations. \textbf{The directions are completely fixed by
$\varphi$} --- no angle remains choosable, only which of the six axes,
and A120 shows that this is a discrete two-fold choice ($2/9$ or
$4/9$), not a free parameter.

The uniqueness has been tested there as well: 200 random rotation axes
at $72^\circ$ give values between $0.036$ and $0.452$, not a single one
hitting $2/9$. For a three-fold axis the result is $2/5$, at $144^\circ$
it is $4/9$.

\section{And the window as well}

An icosahedron has no shape parameter. Its circumradius follows from
the edge:
\[
  |\text{vertex}| = \sqrt{1 + \varphi^2} = 1.9021130, \qquad
  \frac{\text{circumradius}}{\text{edge}} = 0.9510565
  = \frac{\sqrt{\varphi\sqrt5}}{2}.
\]
All size ratios are determined by $\varphi$. What remains is one scale
factor --- and that is exactly the single quantity A085 carries as the
anchor.

\section{The juxtaposition}

\begin{center}
\begin{tabular}{llll}
\toprule
 & cut direction & window shape & window size \\
\midrule
cut-and-project, generic & two free angles & free & free \\
icosahedral & 6 axes, discrete & the icosahedron, shape-fixed & one scale only \\
\bottomrule
\end{tabular}
\end{center}

A continuum of freedoms on one side, a discrete selection plus one
scale on the other.

\section{What this means}

\begin{quote}
Where $\varphi$ is chosen, it has been put in.\\
Where $\varphi$ is found, it fixes direction and window.
\end{quote}

This is not wordplay but the difference between a parameter and a
structure. A construction in which $\varphi$ is set as the projection
direction has a free angle. One in which $\varphi$ occurs as the
eigenvalue of a group has a discrete selection.

And it reverses the touchstone. Instead of only asking \emph{who chose
the direction}, one can ask:

\begin{quote}
\textbf{Is there a structure whose eigenvalues enforce the direction?}
\end{quote}

If an answer comes, the direction is no longer a stipulation.

\chapter{Two conditions instead of three routes}

If $\varphi$ enforces direction and window, Route~2 no longer costs any
free parameters --- as little as Route~1 with a bound radius. The
division into three routes thereby becomes obsolete, but not entirely:
a second condition remains, and it is independent of $\varphi$.

\section{What becomes transferable}

\textbf{The balance of freedoms.} If a structure enforces the
assignment, it is no longer a stipulation, whether one calls it
compactification or projection.
\begin{center}
\begin{tabular}{ll}
Route 1 with free radius & one parameter \\
Route 1 with $\tilde T \cdot m = 1$ & none \\
Route 2 with free direction & several parameters \\
Route 2 with icosahedral axis & none \\
\end{tabular}
\end{center}

The distinction thus does not run between the routes, but between
\emph{enforced} and \emph{stipulated}.

\section{What does not become transferable}

\textbf{The conservation of information.} A projection is by definition
not injective. Even with an enforced direction, $(1,0,0,+1)$ and
$(1,0,0,-1)$ merge into the same image $(1,0,0)$. Twelve winding
vectors become six images, no matter who fixed the direction.

Under compactification this does not happen. The winding number remains
distinguishable --- $+1$ and $-1$ are different states and carry
different quantum numbers, although space has three dimensions.

\section{Where the two nevertheless meet}

Under compactification the information does not vanish; it
\textbf{migrates}. From position into an index. What was a fourth
coordinate before is a winding number afterwards.

A projection with an enforced direction does the same --- \textbf{provided
the preimage pair carries a conserved index}. If it carries none, it is
loss. If it carries one, it is a rebooking, and then the difference
from compactification is merely one of phrasing.

\section{The rule}

\begin{quote}
Not three routes, but two conditions.

\textbf{1. Is the assignment enforced by a structure or stipulated by
the observer?}

\textbf{2. Is the lost positional information caught by a conserved
index?}

If both are met, it does not matter what the procedure is called. If
the first fails, one has parameters. If the second fails, one has loss.
\end{quote}

\section{What this means for the examples}

\begin{center}
\begin{tabular}{lll}
\toprule
 & condition 1 & condition 2 \\
\midrule
Kaluza-Klein, free radius & no & yes \\
cut-and-project, generic & no & no \\
FFGFT, $\tilde T \cdot m = 1$ & yes & yes \\
icosahedral projection & yes & only with an index \\
\bottomrule
\end{tabular}
\end{center}

The last row is the interesting one. It is the case in which $\varphi$
removes the stipulation but the question of the index remains open ---
and that question cannot be answered by geometry, only by stating which
quantity carries the index.

In FFGFT it is the winding number on the mass direction. In a pure
projection construction the corresponding quantity would have to be
named, otherwise the loss is real.

\chapter{The finding that decides the matter}

D4 does not need to be reduced at all. It decomposes by itself.

Of the 24 shortest vectors:
\begin{center}
\begin{tabular}{lcl}
\toprule
 & count & meaning \\
\midrule
with $x_4 = 0$ & \textbf{12} & lie entirely in the first three coordinates \\
with $x_4 \neq 0$ & \textbf{12} & six with $x_4 = +1$, six with $x_4 = -1$ \\
\bottomrule
\end{tabular}
\end{center}

And the twelve with $x_4 = 0$ are not just anything:
\[
  \text{kissing number of D3} = \text{FCC} = 12 .
\]

FCC stands for \emph{face-centered cubic} --- the densest sphere
packing in three dimensions, stacking sequence ABCABC. A sphere in the
interior touches six neighbours in its layer, three above and three
below: $6 + 3 + 3 = 12$. That this is the maximum goes back to the
question of Newton and Gregory and was finally proved in 1953.

Lattice and packing are the same object in two languages: the shortest
vectors of D3 are $(\pm1, \pm1, 0)$ and permutations, exactly twelve,
and they point to the twelve contact points.

\section{Why the pyramid is fundamental in three dimensions}

The twelve neighbour vectors form, as a convex hull, a
\textbf{cuboctahedron} --- eight triangles, six squares, all vertices
at distance $\sqrt2$. This is the neighbourhood figure of
three-dimensional space.

And the same packing has a second shape everyone knows. Placing spheres
into the hollows of a square layer --- the orange stack --- each sphere
in the interior has
\[
  4 \text{ neighbours below} + 4 \text{ in the plane}
  + 4 \text{ above} = 12
\]
Exactly the kissing number. \textbf{The square pyramid is FCC}, merely
stood up differently: the view runs along the cube diagonal instead of
along the edge. Hexagonal layers in the sequence ABCABC and square
layers with hollows describe the same lattice.

The fundamental point is uniqueness. In three dimensions there is
exactly \textbf{one} densest packing --- the Kepler conjecture, shown
by Hales in 1998 and formally verified in 2014. The pyramid is thus not
one building form among several, but the form the density enforces.
This is why copper, gold and aluminium crystallise this way, and why
the fruit stand works.

For this note two things matter in that. First, the 12 is not a
conventional number but provenly maximal in 3D. Second,
three-dimensional space thereby has a distinguished neighbourhood
figure against which one can measure whether a construction computes in
it or not: whoever works spatially with 24 does not work in the
cuboctahedron.
\[
  24 = 12 + 12
\]

\textbf{Half of the D4 kissing number is three-dimensional space. The
other half leads out of it} --- and becomes winding neighbours as soon
as the fourth direction is compact, six in each direction of
circulation.

This is the lattice form of
$\text{T}^4 = \text{T}^3 \times S^1_m$: three spatial circles plus the
fourth direction, and the split $12 + 12$ falls out instead of being
put in.

\chapter{What follows for the invariants}

The invariants of D4 are \textbf{4D quantities}. They are spatial
quantities only if all four directions are space.

\begin{center}
\begin{tabular}{lll}
\toprule
 & D4 (4D) & D3 $=$ FCC (3D) \\
\midrule
kissing number & 24 & 12 \\
packing density & $\pi^2/16 = 0.616850$ & $\pi/(3\sqrt2) = 0.740480$ \\
\bottomrule
\end{tabular}
\end{center}

Whoever computes with 24 as a spatial kissing number thereby implicitly
asserts four spatial dimensions. Whoever has three spatial dimensions
plus one other direction computes spatially with 12, and the other
twelve neighbours are something else --- not neighbours in space, but
windings.

This is no small matter. It affects every quantity formed from the
neighbourhood number.

\section{Which packing density belongs in \texorpdfstring{$K^{-36}$}{K to the -36}?}

From this a question follows for the author's own construction, and it
had to be checked: the bulk exponent of the sector ladder (A270) is
bound to a packing density. If only three of the four directions are
space --- is it then the 4D density or the 3D density?

\begin{center}
\small
\begin{tabular}{lllll}
\toprule
 & inverse density & required exponent
 & distance to $k^*/100$ & to integer \\
\midrule
D4 (4D) & $16/\pi^2 = 1.6211389$ & \textbf{35.9926}
  & $0.098$ = \textbf{0.27\,\%} & $0.0074$ \\
D3 $=$ FCC (3D) & $3\sqrt2/\pi = 1.3504745$ & $22.3836$
  & $13.707$ = \textbf{38.08\,\%} & $0.3836$ \\
\bottomrule
\end{tabular}
\end{center}

And at integer exponent:
\begin{center}
\begin{tabular}{lllr}
D4: & $(75/74)^{36} = 1.621301$ & against $1.621139$ & $+0.010\,\%$ \\
D3: & $(75/74)^{22} = 1.343538$ & against $1.350474$ & $-0.514\,\%$ \\
\end{tabular}
\end{center}
(In the notation $K^{-n}$ with $K = 74/75$; written out as $(75/74)^n$
to keep the sign of the exponent unambiguous.)

\textbf{The result is unambiguous.} The 3D density does not fit. It
requires exponent $22.38$ instead of $36.09$, i.e.\ $38\,\%$ off, and
it additionally lies close to a rounding boundary --- $0.38$ against
$0.007$ for the 4D density. Even as an isolated hit it would be weak.
\textbf{The 4D density belongs in $K^{-36}$.}

\subsection*{Why this is not a contradiction}

The bulk is T$^4$, not T$^3$. The fourth direction contributes to the
packing even though it is not space --- because packing happens where
the torus is, and the torus has four directions.

This is admissible, but it must be said. The mass direction is a
genuine torus direction with integer windings, not a caption. It counts
in the packing because it is geometrically present, and it does not
count in the spatial kissing number because it is not space.

The two statements thus stand side by side without contradicting each
other:
\begin{center}
\begin{tabular}{lcl}
spatial neighbourhood & $\longrightarrow$ & 12 (FCC, the $x_4 = 0$ vectors) \\
packing of the bulk & $\longrightarrow$ & $\pi^2/16$ (all four directions) \\
\end{tabular}
\end{center}

\subsection*{And what follows for comparable constructions}

The same number, two different justifications --- and typically only
one of them is spoken aloud.

In FFGFT, $\pi^2/16$ is the packing density of a four-dimensional torus
of which three directions are space and one carries the mass. That is a
statement that can be justified and is justified.

In a construction that computes on the same lattice but treats all four
directions as spatial, $\pi^2/16$ is the packing density of a
four-dimensional \emph{space}. That is an assertion about the number of
spatial dimensions --- and it is rarely spoken aloud, because the
number appears as a neutral lattice constant.

It is not neutral. Whoever uses it says something about the number of
spatial dimensions, whether they want to or not.

\chapter{Application: constructions with four spatial lattice
directions}

The case that occurs in practice: a construction computes on D4 with 24
offsets, has a step counter for the iteration, and nowhere declares a
time direction. From this it follows necessarily: \textbf{all four
coordinates are space.}

That is an assertion about the world, not merely about the model --- an
object in four spatial dimensions is not a particle in our space, but
an entity in a space that does not exist. Even with a matching node
count it would remain open what that count has to do with a mass.

And the ways out are uncomfortable:
\begin{itemize}
\item One direction becomes time after the fact $\to$ then
  \textbf{three} spatial dimensions remain, and the spatial kissing
  number is 12, not 24. Everything sitting on 24 would have to be
  recomputed.
\item Time is appended as a fifth direction $\to$ then there are four
  spatial dimensions, one too many.
\item It stays at $4+0$ $\to$ then one must say how a four-dimensional
  spatial object acquires a mass in our space.
\end{itemize}

This is the fundamental question for any such construction, and it
precedes the numerical findings: as long as it is open, any correction
to individual quantities addresses only symptoms.

\textbf{It is mostly not asked.} Where time appears neither in the code
nor in the documents nor in the correspondence, the question does not
yet exist for those involved --- and that is understandable, for an
iteration \emph{feels} like time without being one.

\textbf{Asked, not answered.} Whoever builds such a construction knows
its internals; this note reads from outside. There may be answers not
visible from here.

\chapter{The two conditions, applied to string theory}

The rule from the chapter \enquote{Two conditions} can be applied, and
the result is not sweeping. It names exactly one place.

\section{Condition 2 is met exemplarily}

The compact directions are compact. Nothing is projected, nothing
merges. Winding numbers, Kaluza-Klein momentum, winding sectors, Hodge
numbers, intersection numbers --- all conserved indices. The positional
information of the six rolled-up directions migrates completely into
quantum numbers.

This is the strength of the approach and should not be passed over. In
the language of the rule: the information is rebooked, not lost.

\section{Condition 1 is met by half}

The dividing line runs through the middle of the construction.

\textbf{Enforced:}
\begin{center}
\begin{tabular}{l}
$D = 10$ from anomaly freedom and the critical dimension \\
$N = 1$ SUSY in 4D demands SU(3) holonomy, hence Calabi-Yau \\
from that: $10 - 4 = 6$ compact directions \\
\end{tabular}
\end{center}

The \textbf{number} of rolled-up directions is not choosable. That is a
genuine enforcement, and it is stronger than anything FFGFT can adduce
for its four circles --- there the choice of T$^4$ is marked as
[STIPULATION].

\textbf{Stipulated:}
\begin{center}
\begin{tabular}{ll}
which Calabi-Yau manifold & ($\sim 10^5$ known, total number unknown) \\
flux quantum numbers & ($\sim 10^{500}$ combinations) \\
K\"ahler and complex structure moduli & (continuum, then stabilised) \\
standard embedding or not & \\
\end{tabular}
\end{center}

The number of directions is enforced. The \textbf{geometry within} is
not.

\section{What the landscape thereby is}

\begin{quote}
The landscape is exactly the unmet part of condition 1.
\end{quote}

Not an error and not a scandal, but the place where the structure stops
enforcing. A precise statement of location instead of an accusation.

And it explains why the problem is so stubborn: it cannot be fixed by
more computation. Where a structure does not enforce, no computation
can make it do so. An additional condition from outside is needed ---
and that is exactly what has been sought for decades.

\section{The asymmetry that connects the two}

What is compactified are \textbf{spatial} directions. Time stays
outside. By A060 it follows: no dynamical scale in the state space. All
scales must come from elsewhere, and in this construction they come
from the moduli.

But moduli are exactly what condition 1 does not cover.
\begin{center}
\begin{tabular}{l}
time outside \ $\to$ \ no scale in the structure \\
\phantom{time outside} $\to$ \ scales from moduli \\
\phantom{time outside} $\to$ \ moduli are the unenforced part \\
\phantom{time outside} $\to$ \ landscape \\
\end{tabular}
\end{center}

The landscape arises where the scales must come from. And the scales
must come from there because time stands outside. It is a single choice
that produces both: roll up six spatial directions and leave time open.

\section{What this does not say}

It does not say that string theory is wrong. It meets condition~2
better than anything considered here, and its dimension count is more
compellingly justified than the four circles of FFGFT.

Nor does it say that the alternative would be cheaper. A rolled-up time
costs causality objections, Pauli's theorem and the thermal reading ---
and it is nowhere accepted so far.

What it says is exclusively this: \textbf{the landscape and the
necessity of attached scales are not two problems but one.} Whoever
wants to solve one must start at the other.

\chapter{Verification script}

\begin{verbatim}
import itertools, math, collections

D4 = [v for v in itertools.product([-1,0,1], repeat=4)
      if sum(x*x for x in v) == 2]
D3 = [v for v in itertools.product([-1,0,1], repeat=3)
      if sum(x*x for x in v) == 2]

print(len(D4))                                 # 24
print(sum(1 for v in D4 if v[3] == 0))         # 12
print(sum(1 for v in D4 if v[3] != 0))         # 12
print(len(D3))                                 # 12  (FCC)
print(len(set(v[:3] for v in D4)))             # 18  (projection)
print(collections.Counter(
    v[3] for v in D4 if v[3] != 0))            # {1: 6, -1: 6}

print(math.pi**2 / 16)                         # 0.616850  D4
print(math.pi / (3 * 2**0.5))                  # 0.740480  D3
\end{verbatim}

Extended version with all numbers of this document:
\texttt{python/\allowbreak Dok311\_\allowbreak Skripte/\allowbreak
ffgft\_\allowbreak 311\_\allowbreak vier\_\allowbreak auf\_\allowbreak
drei.py}.

\chapter{References}

\textbf{Internal documents:} A060 (time as scale), A085 (anchor), A120
($\theta = 2/9$, five-fold rotation), A270 (bulk exponent 36, P35),
Doc.~275 (scale recursion depth $k^*$), Doc.~310 (resonance geometry as
reference, concert-pitch anchor), Register 190 (R67: uncertainty of the
factor 100 in the denominator of $k^*/100$).

\medskip
\noindent\textbf{On status:} working note for the author's own use. The
chapters on the routes, on the decomposition finding $24 = 12+12$ and
on the packing density are computed; the chapter
\enquote{Application} is a question and not an assertion.

\end{document}
